\documentclass{article}
\usepackage{amsmath}
\usepackage{amssymb}
\usepackage{amsthm}

\usepackage{color}
\newcommand{\red}[1]{ {\color{red}{#1}} }
\newcommand{\blue}[1]{ {\color{blue}{#1}} }

\newcommand{\N}{\mathbb{N}}
\newcommand{\Z}{\mathbb{Z}}
\newcommand{\Q}{\mathbb{Q}}
\newcommand{\R}{\mathbb{R}}
\newcommand{\C}{\mathbb{C}}

\theoremstyle{plain}
    \newtheorem{definition}{Definition}
    \newtheorem{example}{Example}
    \newtheorem{proposition}{Proposition}

\title{Writeup}
\author{Caleb and Jeremy}
\date{September 2026}

\begin{document}
\maketitle

\section{Language}


\subsection{Introduction}

The goal of this section is to characterize when the function defined by
\[ 
    \Z/n\Z \xrightarrow[x \mapsto ax+b]{} Z/n\Z
\]
is invertible.
That will give conditions for an \textit{affine cipher} to be a valid encryption method.

To do this we will need some basic number theory, and the following correspondence
between English letters and integers:
\begin{center}
\begin{tabular}{c|c|c|c|c|c}
    A & B & C & $\cdots$ & Y & Z  \\ \hline
    0 & 1 & 2 & $\cdots$ & 24 & 25
\end{tabular}
\end{center}


\subsection{Basic Number Theory}

The symbol $\Z$ denotes the set
\[
    \{ \dots, -2, -1, 0, 1, 2, \dots\}
\]
of \textbf{integers} or \textbf{whole numbers}.

\begin{definition}
    Let $n$ be a positive integer.
    Define $\Z/n\Z$ to be the set 
    \[
        \{ [0], [1], \dots, [n-1] \}.
    \]
    Sometimes, instead of writing $[k]$ we will write $k\pmod{n}$
    or simply $k$ if there is no confusion.
\end{definition}

In some sense, we can think of $\Z/n\Z$ as a collection of all possible 
remainders upon division by $n$.
This idea is illustrated in the following fact.

\begin{proposition}[Division]
    Let $a,n \in \Z$ be two positive integers.
    There is exactly one pair $q,r\in\Z$ of positive integers satisfying 
    \begin{enumerate}
        \item $a = qn+r$,
        \item $0\leq r < n$.
    \end{enumerate}
\end{proposition}

\blue{
\begin{proof}
    To do.
\end{proof}
}

In the case when $r=0$, we say $n|a$, pronounces \textbf{$n$ divides $a$}.\footnote{
    The symbol $|$ has an alternative pronunciation: ``guzinta.''
As in, ``$n$ guzinta $a$,'' or, ``$n$ goes into $a$.''}
We refer to $r$ as the remainder upon division by $n$, or 
the \textbf{residue of $a$ modulo $n$}.
We express this by writing $a \equiv r \pmod{n}$.

Two integers may have the same residue modulo $n$.
If $a \equiv r \pmod{n}$, then $a+n \equiv r \pmod{n}$ since 
\[ 
    a+n = (q+1)n + r.
\] 
A similar calculation shows that
\[
    r \equiv a+2n \equiv a+3n \equiv \cdots \pmod{n}.
\]
Thus we think of $r$ as ``representing'' the family $a+kn$ of integers, where $k$
ranges across all elements of $\Z$.








\subsection{Caesar Ciphers and Affine Ciphers}

We begin by considering 



\end{document}